\section{Introduction}

The incidence matrix of a block design records which points belong to each
block. Its integral column lattice describes the integer point vectors that
can be assembled from blocks. Every such vector has coordinate sum divisible
by the block size, but this necessary condition need not be sufficient.
The quotient measuring this failure is the object of this paper.

Let $D=(P,\mathcal B)$ be a nonempty, possibly non-simple
$t$-$(v,k,\lambda_t)$ design with $2\le t\le k<v$. Write $A$ for its
point-by-block incidence matrix and set
\[
L=\im_\ZZ(A),\qquad
\Hk=\Big\{z\in\ZZ^P:\textstyle\sum_{x\in P}z_x\equiv0\pmod k\Big\}.
\]
Since every column has weight $k$, one has $L\le\Hk$.

\begin{definition}
The \emph{selected higher-star defect} of $D$ is the quotient
$\Delt(D)=\Hk/L$.
\end{definition}

The word \emph{selected} refers to the actual blocks of $D$, rather than the
complete collection of $k$-subsets. Theorem~\ref{thm:hs} shows that the defect
is finite. It is the kernel of the coordinate-sum map
$\coker(A)\to\ZZ/k$, so it removes the cyclic factor forced by the common
block size. We study its bounds, its dependence on the chosen design and
its behavior under local changes of blocks.

\subsection{Bounds from the design parameters}

Classical star identities give an annihilator obtained from the successive
incidence numbers. After this elementary reduction in Section~\ref{sec:stars},
we replace $A$ by a pointed matrix presenting the defect itself. Its Gram
matrix depends only on the ordinary $2$-design parameters. The dual defect
then sits as an isotropic subgroup of the associated finite bilinear group
(Proposition~\ref{prop:gp}).

This presentation gives two useful local calculations. At odd primes of
valuation one, an explicit determinant class yields a Witt bound that is
strictly stronger than square divisibility on an infinite family of
$S(2,5,v)$ (Theorem~\ref{thm:wi}). When the usual zero-sum decomposition is
unavailable, a rank-three matrix contains the remaining local information
(Theorem~\ref{thm:xd}). For binary subgroup orders, the answer is simpler:
a general bilinear lemma shows that isotropy alone gives no improvement on
the square-divisibility bound (Corollary~\ref{thm:xi}).

\subsection{Variation, local changes and realization}

The parameters do not determine the defect. Two Steiner triple systems on
$15$ points have the same parameter Gram matrix but different selected
quotients (Theorem~\ref{thm:sep}). When a symmetric design has order of
$p$-adic valuation two, however, one modular rank determines the local group
(Theorem~\ref{thm:v2}). Section~\ref{sec:parameters} treats these two phenomena.

A Pasch switch replaces four triples by four others on the same six points.
Theorem~\ref{thm:pt} shows that it preserves every odd-primary part of the
defect and changes the total order by a factor in $\{1/2,1,2\}$.
The proof follows the two local column lattices after the unchanged blocks
have been adjoined.

The realization result has an infinite scope: every fixed elementary binary
or ternary group occurs for all sufficiently large orders in the respective
classical Steiner triple-system families (Theorem~\ref{thm:er}). Appendix~\ref{sec:app-er}
proves the required asymptotic comparison of published exact-rank counts.
The geometric applications and worked examples show both vanishing and
nonzero defects.

\subsection{Relation to earlier work}

The higher-star identity comes from Wilson's incidence theory
\cite{Wilson1982,KiermaierWassermann}. The universal additive group and the
$t=2$ point-difference identity are treated by Caggegi--Falcone--Pavone
\cite{CFP2017}. Later developments concern additive embeddings of designs
\cite{BurattiNakic,BMN2025}; the point-injectivity questions in that theory
are outside the scope of this paper.

The complete-subset incidence matrix has a classical diagonal form
\cite{Wilson}, and integral design modules are studied in
\cite{GraverJurkat,GrahamLiLi,Keevash}. Here the column lattice belongs to a
particular design. We use the standard theory of finite bilinear forms and
lattice saturation \cite{Wall1963,Nikulin}; applying Witt theory to designs
also has classical precedents \cite[Sections 5--6]{WilsonNotes}.

The deductions developed below are the explicit pointed Gram calculation,
its prime-local consequences, the Pasch quotient law and the fixed-defect
asymptotic comparison. The star bound and binary order lemma are elementary
supporting facts. Geometric Smith forms, finite rank spectra and enumeration
theorems are imported where they are used; Sin's survey \cite{Sin2014}
provides broader context for incidence Smith forms. These results concern
the image lattice and bilinear isotropy. They neither classify relation
kernels nor settle general realization or quadratic-refinement questions.

\subsection{Notation}

For $1\le j\le t$, let $\lambda_j$ be the number of blocks containing a
fixed $j$-set. We use
\[
r=\lambda_1,\qquad \lambda=\lambda_2,\qquad
\alpha=r-\lambda,\qquad n=v-1,\qquad b=|\mathcal B|.
\]
For a finite abelian group $F$, write $F_p$ for its $p$-primary part,
$\ell_p(F)=\log_p|F_p|$ and $\mu(F)$ for its minimum number of generators.
The symbols $\ZZ_p$ and $\FF_p$ denote the $p$-adic integers and the field
of $p$ elements. Unless stated otherwise, a discriminant pairing is
symmetric and bilinear with values in $\QQ/\ZZ$.
